\section{Primitive characters and the critical strip}
\label{sec:primitive-reduction}

This section isolates a classical part of the argument which can be written
without any of the automorphic constructions in the announced proof.  It does
not establish the required nonvanishing in the right half-plane.  It proves,
with the analytic details made explicit, why that assertion would imply the
usual statement about zeros of every Dirichlet $L$-function in the critical
strip.

\subsection{Primitive Gauss sums}

Let $q$ be a positive integer.  A Dirichlet character modulo $q$ is a
homomorphism from $(\mathbb Z/q\mathbb Z)^\times$ to $\mathbb C^\times$,
extended periodically to the integers by zero on integers not relatively
prime to $q$.  For $q=1$ the unique character has value $1$ at every integer.
A character modulo $q$ is \emph{primitive} if it is not induced from a
character modulo a proper divisor of $q$.  Here induction means that the
values agree at all integers relatively prime to $q$.  The conductor of a
character is the modulus of its inducing primitive character.

Every character modulo $N$ is induced from a unique primitive character
modulo a divisor $q$ of $N$.  To see this explicitly, write
$N=\prod_{p\mid N}p^{e_p}$.  The Chinese remainder theorem identifies
its unit group with the product of the unit groups modulo $p^{e_p}$,
and the character is uniquely a product of characters of those factors.
For each $p$ choose the smallest $f_p\in\{0,\ldots,e_p\}$ through which
the corresponding character factors under reduction modulo $p^{f_p}$;
the unit group for $p^0=1$ is trivial.  Reduction maps of these unit groups
are surjective, so the factored characters are unique.  Their product
has modulus $q=\prod p^{f_p}$, induces the original character, and is
primitive by minimality at each prime.  Any primitive inducing modulus
must have exactly these exponents: an exponent below $f_p$ is impossible,
and one above $f_p$ gives an imprimitive character.  This proves the
assertion and the uniqueness.

Every character value at a unit is a root of unity, because the group of
units is finite.  Thus $|\psi(n)|\leq 1$ and
$\overline{\psi(n)}=\psi(n)^{-1}$ for units $n$.
For a primitive character $\psi$ modulo $q>1$, put
\[
 a=\begin{cases}0,&\psi(-1)=1,\\1,&\psi(-1)=-1,\end{cases}
 \qquad
 \tau(\psi)=\sum_{r\bmod q}\psi(r)e^{2\pi i r/q}.
\]
The two possibilities for $\psi(-1)$ follow from $(-1)^2=1$.

\begin{lemma}\label{prim:gauss}
For every integer $m$,
\[
 \sum_{r\bmod q}\psi(r)e^{2\pi i mr/q}
   =\overline{\psi(m)}\tau(\psi),
 \qquad |\tau(\psi)|=\sqrt q.
\]
\end{lemma}

\begin{proof}
If $m$ is relatively prime to $q$, multiplication by $m$ permutes the
residue classes modulo $q$.  Substituting $u=mr$ gives the first identity.

Suppose now that $d=(m,q)>1$, and put $q'=q/d$.  Reduction of units modulo
$q$ onto the units modulo $q'$ is surjective.  Indeed, a residue class
coprime to $q'$ can be lifted to an integer coprime also to every prime
dividing $q$ but not $q'$, by the Chinese remainder theorem.  If $\psi$
were trivial on the kernel of this reduction, it would factor through
the units modulo $q'$, contradicting primitivity.  There is therefore a
unit $u$ modulo $q$ such that $u\equiv 1\pmod {q'}$ and $\psi(u)\neq1$.
Since $m(u-1)$ is divisible by $q$, substitution $r\mapsto ur$ in the sum
shows that the sum equals $\psi(u)$ times itself.  It is zero.  The right
side is also zero because $\psi(m)=0$.

Write $S(m)$ for the sum in the first identity.  Finite Fourier
orthogonality gives
\begin{align*}
 \sum_{m\bmod q}|S(m)|^2
 &=\sum_{r,t\bmod q}\psi(r)\overline{\psi(t)}
       \sum_{m\bmod q}e^{2\pi i m(r-t)/q}\\
 &=q\sum_{r\bmod q}|\psi(r)|^2=q\varphi(q),
\end{align*}
where $\varphi(q)$ is the number of units modulo $q$.
On the other hand, the identity already proved makes the left side
$|\tau(\psi)|^2\varphi(q)$.  Division by the positive integer
$\varphi(q)$ proves the second assertion.
\end{proof}

\subsection{The primitive functional equation}

For a positive real number $x$, complex powers mean
$x^s=\exp(s\log x)$ with the real logarithm.
For $\operatorname{Re}s>1$, define
\[
 L(s,\psi)=\sum_{n=1}^{\infty}\frac{\psi(n)}{n^s}.
\]
The series converges absolutely and locally uniformly in that half-plane.
For primitive $\psi$ modulo $q>1$, define its completion there by
\[
 \Lambda(s,\psi)
 =\left(\frac q\pi\right)^{(s+a)/2}
     \Gamma\left(\frac{s+a}{2}\right)L(s,\psi).
\]
Here $\Gamma$ is Euler's gamma function. Lemma~\ref{tools:gamma}
proves its nonvanishing and the entire continuation of its reciprocal.
We use Poisson summation in the normalization of
Lemma~\ref{tools:poisson} and equation~\eqref{tools:poisson-progression}.

\begin{proposition}\label{prim:functional-equation}
If $\psi$ is primitive modulo $q>1$, then $\Lambda(s,\psi)$ extends to an
entire function and
\[
 \Lambda(s,\psi)
    =\varepsilon_\psi\Lambda(1-s,\overline\psi),
 \qquad
 \varepsilon_\psi=i^{-a}\frac{\tau(\psi)}{\sqrt q},
 \qquad |\varepsilon_\psi|=1.
\]
The unique primitive principal character has conductor $1$.  Its
$L$-function is $\zeta(s)$, and
\[
 \Lambda_0(s)=\pi^{-s/2}\Gamma(s/2)\zeta(s)
\]
has a meromorphic continuation with only simple poles at $0$ and $1$,
and satisfies $\Lambda_0(s)=\Lambda_0(1-s)$.
\end{proposition}

\begin{proof}
First assume $q>1$ and put
\[
 \theta_\psi(t)=\sum_{n\in\mathbb Z}
      \psi(n)n^a e^{-\pi n^2t/q},\qquad t>0.
\]
When $a=0$ the factor $n^a$ is interpreted as $1$; the term $n=0$ is
still zero because $\psi(0)=0$.  The series and all its derivatives
converge uniformly on compact subintervals of $(0,\infty)$.

Applying Poisson summation in each residue class modulo $q$ yields
\[
 \sum_{n\in\mathbb Z}\psi(n)f(n)
   =\frac1q\sum_{m\in\mathbb Z}\widehat f(m/q)
       \sum_{r\bmod q}\psi(r)e^{2\pi i mr/q}
   =\frac{\tau(\psi)}q\sum_{m\in\mathbb Z}
       \overline{\psi(m)}\widehat f(m/q).
\]
Use $f(x)=x^a e^{-\pi tx^2/q}$.  The Gaussian Fourier transform and its
first derivative give, for $a=0,1$,
\[
 \widehat f(y)=(-i)^a
      (q/t)^{a+1/2}y^a e^{-\pi qy^2/t}.
\]
Consequently
\begin{equation}\label{prim:theta-transform}
 \theta_\psi(t)
  =\varepsilon_\psi t^{-a-1/2}
       \theta_{\overline\psi}(1/t).
\end{equation}
Lemma~\ref{prim:gauss} proves $|\varepsilon_\psi|=1$.

The theta series is exponentially small as $t\to\infty$.
Equation~\eqref{prim:theta-transform} shows that as $t\to0^+$ it is
bounded by a fixed power of $t^{-1}$ times $e^{-c/t}$, for some $c>0$.
It follows that
\[
 F_\psi(s)=\frac12\int_0^\infty
           \theta_\psi(t)t^{(s+a)/2-1}\,dt
\]
converges absolutely and locally uniformly for every $s\in\mathbb C$.
The same is true after inserting any fixed power of $\log t$, so
differentiation under the integral proves that $F_\psi$ is entire.

For $\operatorname{Re}s>1$, absolute convergence justifies termwise
integration.  Indeed, after taking absolute values, the integral of the
$n$th term is a constant depending on $s$ times $|n|^{-\operatorname{Re}s}$.
Since $\psi(-n)(-n)^a=\psi(n)n^a$, the factor $1/2$ removes the duplicate
negative terms.  Substitution $u=\pi n^2t/q$ therefore gives
\[
 F_\psi(s)
 =\Gamma\left(\frac{s+a}{2}\right)
      \left(\frac q\pi\right)^{(s+a)/2}
      \sum_{n\geq1}\frac{\psi(n)}{n^s}
 =\Lambda(s,\psi).
\]
Finally insert~\eqref{prim:theta-transform} in the integral and make the
change of variable $u=1/t$.  Absolute convergence at both endpoints
justifies this substitution and yields
\[
 F_\psi(s)=\frac{\varepsilon_\psi}{2}
       \int_0^\infty\theta_{\overline\psi}(u)
           u^{(1-s+a)/2-1}\,du
   =\varepsilon_\psi F_{\overline\psi}(1-s).
\]
This proves the functional equation and the continuation for $q>1$.
Dividing the entire completion by the displayed gamma and power factors
also continues $L(s,\psi)$; equivalently, multiplication by the entire
function $1/\Gamma((s+a)/2)$ shows that $L(s,\psi)$ is entire.

For the principal character of conductor $1$, let
$\theta(t)=\sum_{n\in\mathbb Z}e^{-\pi n^2t}$ and $A(t)=\theta(t)-1$.
Poisson summation gives $\theta(t)=t^{-1/2}\theta(1/t)$.
For $\operatorname{Re}s>1$, termwise integration gives
$\Lambda_0(s)=\frac12\int_0^\infty A(t)t^{s/2-1}\,dt$.
Split the integral at $1$ and transform the portion from $0$ to $1$.
The result is
\begin{equation}\label{prim:zeta-continuation}
 \Lambda_0(s)=\frac1{s-1}-\frac1s
  +\frac12\int_1^\infty A(t)
       \bigl(t^{s/2-1}+t^{(1-s)/2-1}\bigr)\,dt.
\end{equation}
The integral is entire by exponential decay, and the right side is
unchanged by $s\mapsto1-s$.  It has exactly the asserted poles, with
residues $1$ at $1$ and $-1$ at $0$.
\end{proof}

\subsection{Changing the modulus}

Suppose $\chi$ modulo $N$ is induced by a primitive character $\psi$
modulo $q$, where $q\mid N$.  Define the entire function
\[
 E_{N,q}(s,\psi)
   =\prod_{\substack{p\mid N\\p\nmid q}}
         (1-\psi(p)p^{-s}),
\]
where $p$ runs over primes and an empty product is $1$.

\begin{lemma}\label{prim:change-level}
As an identity of meromorphic functions,
\[
 L(s,\chi)=L(s,\psi)E_{N,q}(s,\psi).
\]
Moreover, $E_{N,q}(s,\psi)\neq0$ whenever
$\operatorname{Re}s>0$.  Thus $L(s,\chi)$ and $L(s,\psi)$ have the same
zeros, with the same multiplicities, in that open half-plane, away from
any pole.
\end{lemma}

\begin{proof}
In $\operatorname{Re}s>1$, absolute convergence and unique prime
factorization give the Euler products
\[
 L(s,\chi)=\prod_p(1-\chi(p)p^{-s})^{-1},
 \qquad
 L(s,\psi)=\prod_p(1-\psi(p)p^{-s})^{-1}.
\]
For primes $p\nmid N$ the character values agree.  For primes dividing
$q$ both values vanish.  For $p\mid N$, $p\nmid q$, the factor of
$L(s,\chi)$ is $1$, whereas the factor of $L(s,\psi)$ is
$(1-\psi(p)p^{-s})^{-1}$.  Cancelling these factors proves the identity
in that half-plane.  The identity supplies the meromorphic continuation
of $L(s,\chi)$ and remains valid everywhere by uniqueness of analytic
continuation.

If $\sigma=\operatorname{Re}s>0$, then
\[
 |\psi(p)p^{-s}|\leq p^{-\sigma}<1,
 \qquad
 |1-\psi(p)p^{-s}|\geq1-p^{-\sigma}>0.
\]
Every factor, and hence the finite product, is nonzero.  Multiplication
by a holomorphic function which is nonzero at a point preserves the
order of a zero there, proving the multiplicity statement.
\end{proof}

For the primitive principal character, we use the conventions
$a=0$ and $\Lambda(s,\psi)=\Lambda_0(s)$ from
Proposition~\ref{prim:functional-equation}.

The restriction $\operatorname{Re}s>0$ in this lemma is essential.
For example, for the principal character modulo a prime $p$,
$L(s,\chi)=\zeta(s)(1-p^{-s})$ as meromorphic functions.
The extra factor vanishes at $s=2\pi i k/\log p$ for $k\in\mathbb Z$.
In particular, $s=0$ is an additional zero since $\zeta(0)=-1/2$.
This value follows directly from~\eqref{prim:zeta-continuation}:
$\Lambda_0(s)\sim-1/s$ and $\Gamma(s/2)\sim2/s$ at zero.
Thus imprimitive characters require care on the boundary of the critical
strip.

\subsection{The precise reduction}

For clarity, state the needed half-plane assertion explicitly:
\begin{quote}
For every positive integer $N$, every Dirichlet character $\chi$ modulo
$N$, and every $s\in\mathbb C$ with $\operatorname{Re}s>1/2$, the value
$L(s,\chi)$ is nonzero whenever $\chi$ is nonprincipal or $s\neq1$.
\end{quote}
At $s=1$ a principal $L$-function has a pole, so that exception is part of
the statement, not a possible zero.  The value of a meromorphic function
at its pole is not used in the argument.

\begin{theorem}[Critical-strip reduction]\label{prim:strip-reduction}
Assume the half-plane assertion above.  Then for every Dirichlet
character $\chi$ modulo a positive integer $N$,
\[
 0<\operatorname{Re}s<1,\qquad L(s,\chi)=0
 \quad\Longrightarrow\quad
 \operatorname{Re}s=\frac12.
\]
In fact, it suffices to assume nonvanishing for primitive characters at
points satisfying $1/2<\operatorname{Re}s<1$.
\end{theorem}

\begin{proof}
Let $\psi$ be the primitive character inducing $\chi$.
Lemma~\ref{prim:change-level} shows that a zero $s$ of $L(s,\chi)$ in
the strip is a zero of $L(s,\psi)$.  If
$\operatorname{Re}s>1/2$, the stated nonvanishing already gives a
contradiction.

Suppose instead that $\operatorname{Re}s<1/2$.  The gamma factor and
the power of $q/\pi$ in the completion are finite and nonzero at $s$,
because $\operatorname{Re}((s+a)/2)>0$.  Therefore
$\Lambda(s,\psi)=0$.  Proposition~\ref{prim:functional-equation} and
$\varepsilon_\psi\neq0$ imply
$\Lambda(1-s,\overline\psi)=0$.  For the character of conductor $1$
the same inference uses~\eqref{prim:zeta-continuation}; neither $s$ nor
$1-s$ is a pole.  The conjugate character is primitive: any factorization
of $\overline\psi$ through a proper divisor, when conjugated, would give
one for $\psi$.  Its parity is the same as that of $\psi$.

At $1-s$ the gamma and power factors are again finite and nonzero,
since $0<\operatorname{Re}(1-s)<1$.  Hence
$L(1-s,\overline\psi)=0$.  But
$1/2<\operatorname{Re}(1-s)<1$, contradicting the assumed primitive
nonvanishing.  The only remaining possibility is
$\operatorname{Re}s=1/2$.
\end{proof}

The character changes to its conjugate in this reflection.  For a
nonreal character it is generally incorrect to replace
$\overline\psi$ by $\psi$ in the displayed functional equation.
No bound on the conductor, and no restriction to real or even
characters, is present in the theorem.

The reduction just proved is the classical part of the passage from
half-plane nonvanishing to the conclusion of the announced
result~\cite{AlpogeSource}. The half-plane estimate itself requires
further arithmetic input.
